\subsection{Short-terminal fourth moments with an arbitrary receiving arm}
\label{proof:Q5-C117}

All operators, expectations, and paths use the original finite
probability space of Lemma~\ref{lem:reference-common:receiving}.
Let \(a\ge3\) be the receiving-arm length, and put
\(h_2=Tp=q-Kg\). The receiving arm and the terminal \(b\)-arm
point outward; every connector points from terminal to receiver.
For \(h_2\), the internal vertex of the two-edge arm points to
the terminal and to its leaf. These directions depend only on the
arm lengths. Let \(E_o,E_t\) denote arbitrary actual equal-pair
products at the receiving and terminal vertices, respectively.

\begin{proposition}[The three receiving families]
\label{prop:q117-unified:families}
For every legal original model and all the indicated integers,
\begin{align}
 \E[E_of_a\{T^L(E_tqf_b)\}^4]
   &\le1-d/3000,& b\ge3,\quad L\ge1,
   \label{eq:q117-unified:shortone}\\
 \E[E_of_a\{T^L(E_tqf_2)\}^4]
   &\le1-3d/800,& L\ge6,
   \label{eq:q117-unified:shortone-two-tail}\\
 \E[E_of_a\{T^L(E_th_2f_b)\}^4]
   &\le1-d/400,& b\ge3,\quad L\ge2.
   \label{eq:q117-unified:shorttwo}
\end{align}
\end{proposition}

The proof uses nine finite positive envelopes. Their actual means,
variances, and signed receiving covariances feed the same fourth
identity. Tables of finite parameters then cover the short integer
prefixes and the unbounded connector tails. The high domains use
direct positive path comparisons.

\subsubsection{A finite receiving-symbol bound}
\label{sec:q117-unified:symbols}

The reversed positive path convolution gives
\fitdisplay{
 S_n=1+G_n(K)g
   =f_n+\sum_{j=3}^n\beta_jf_{n-j},\qquad
 \tau\le\beta_j\le d\quad(j\ge3).
 \label{eq:Q117:conv}
}{Q117.1.1}
Consequently the weak receiving upper is
\fitdisplay{
 0\le f_a\le C_{a,0}=S_a-\beta_a\le2,\qquad
 \E C_{a,0}=1-\beta_a-\beta_{a+1}.
 \label{eq:Q117:receiver}
}{Q117.1.2}
The upper bound also follows directly from
\(C_{a,0}=AS_{a-1}\), \(S_{a-1}\ge0\), and
\(\E S_{a-1}=1-\beta_a\), using the original kernel bound
\(A\le2\).
For \(a\ge4\), the next actual positive term permits
\(C_{a,\beta_{a-1}}=S_a-\beta_a-\beta_{a-1}p\ge f_a\).
These are the common weights of
Lemma~\ref{lem:reference-common:receiving}.
The original spectral measure \(\nu\) has support
\([0,\theta]\), mass \(d\), and
\(\int x^j\,d\nu=\delta_{j+1}\). For \(d>0\),
\(d\le v:=\delta_2/d\le\theta\); the lower bound follows
from \(\E(Kg)=-d\) and Cauchy's inequality.
This constraint concerns aggregate moments, rather than each
spectral frequency.

\begin{lemma}[All finite short-one receiving symbols]
\label{lem:q117-unified:shortone-symbol}
For \(b\ge3\), put \(N=b-1\ge2\) and \(1\le k\le6/5\).
For \(b=2\), put \(N=k=1\). For every integer \(L\ge1\),
let \(P_L(z)=G_L(z)+kz^{L+1}G_N(-z)\). Then for every
\(a\ge3\), \(0\le x\le\theta\), and \(z=i\sqrt x\),
\begin{equation}\label{eq:q117-unified:shortone-symbol}
 \operatorname{Re}(G_a(z)P_L(z))
 \ge\phi(x):=1-3x+2x^2-x^3
 \ge1451/32768>1/25.
\end{equation}
\end{lemma}
\begin{proof}
Put \(c=1-x\). Write \(G_a=u_a+zv_a\); finite alternating
sums give \(c\le u_a,v_a\le1\). For \(N\ge2\), write
\(G_N(-z)=\alpha_0-z\beta_0\). The same finite sums give
\[
 c\le\alpha_0,\beta_0\le1,\qquad
 0\le\alpha_0-c\beta_0\le x,\qquad
 c\le c\alpha_0+x\beta_0\le1.
\]
Indeed for \(N=2h\), both coefficients are
\(a_h=(1-(-x)^h)/(1+x)\), so their first difference is
\(xa_h\). For \(N=2h+1\), it is
\((x+(-x)^h)/(1+x)\), lying in \([0,x]\) for every
\(h\ge1\). The other inequalities are the corresponding finite
alternating endpoints. For \(N=1\), use separately
\(\alpha_0=1,\beta_0=0\).

Write \(P_L=r_L+zs_L\). For \(N\ge2\), the first two rows are
\((r_1,s_1)=(1-kx\alpha_0,kx\beta_0)\) and
\((r_2,s_2)=(1-kx^2\beta_0,1-kx\alpha_0)\).
Both coordinates are nonnegative. For \(L\ge3\), each unperturbed
coordinate is at least \(c\), and each possible negative
correction has magnitude at most \(kx^2\). Hence
\(r_L,s_L\ge c-kx^2\ge1013/2560>0\).
The same argument works for \(N=k=1\), allowing \(s_1=0\).
It follows that
\[
 \operatorname{Re}(G_aP_L)=u_ar_L-xv_as_L
 \ge cr_L-xs_L=\operatorname{Re}((c+z)P_L).
\]

Put \(B_0=(1-2x)/(1+x)\), \(A_0=(2-x)/(1+x)\), and
\(D_0=A_0-c=(1-x+x^2)/(1+x)>0\).
Then \(D_0+A_0x=1\), \(kx\le1\),
and \(B_0-D_0x^2=\phi(x)\). Finite geometric division gives
\[
 \operatorname{Re}((c+z)G_{2m})=B_0(1-(-x)^m),\qquad
 \operatorname{Re}((c+z)G_{2m+1})
       =B_0+(-1)^mA_0x^{m+1}.
\]
When \(N\ge2\), the \(L=1\) value is at least \(c-kx\),
whose excess over \(\phi\) is at least
\(x(4/5-2x+x^2)\ge(781/5120)x\).
The \(L=2\) value is
\(1-2x+kx^2(\alpha_0-c\beta_0)\ge1-2x\ge\phi\).
For even \(L=2m\ge4\), the correction is
\(-k(-1)^mx^{m+1}(\alpha_0-c\beta_0)\).
Its negative part costs at most \(kx^4\le x^3\), paid by
the unperturbed lower bound \(c(1-2x)\).
For odd \(L=2m+1\ge3\), put
\(C_0=c\alpha_0+x\beta_0\in[c,1]\).
The exact value is \(B_0+(-1)^m(A_0-kC_0)x^{m+1}\).
For odd \(m\), this is at least \(B_0-D_0x^{m+1}\ge\phi\).
For even \(m\ge2\), a negative correction can occur only if
\(k>A_0\); then \((k-A_0)x\le D_0\), so it costs at most
\(D_0x^2\).

For \(N=k=1\), the even correction is
\(-(-1)^mx^{m+1}\). At \(m=1\) the result is
\((1-x)^2\ge\phi\); at \(m\ge2\) its negative part is at
most \(x^3\). The odd value, including \(L=1\), is exactly
\(B_0+(-1)^mD_0x^{m+1}\), again at least \(\phi\).
Finally \(\phi'=-3+4x-3x^2<0\), and
\(\phi(\theta)=1451/32768\).
\end{proof}

\paragraph{Q117.4 The wider coefficient-nine entrance}
\label{part:Q117.4}
This entrance also applies when the receiving upper is not one of
the \(a\ge3\) weights of the proposition. Suppose \(C,Z\ge0\),
\(u_0=\E Z\ge0\), and
\fitdisplay{
 C\le p+(6/5)\sqrt d,\qquad
 Z\le(1-5d/8)q-d/2+(13/12)\sqrt d,\qquad
 u_0\le1-3d/2.
}{Q117.4.1}
Then on \(0\le d\le1/3\),
\fitdisplay{
 C[(Z+u_0)^2+2u_0^2]\le9.
 \label{eq:q117-unified:weak-nine}
}{Q117.4.2}
Indeed \(t=\sqrt d\) gives exactly case \(w\) of
Lemma~\ref{lem:reference-coefficient-nine}, with the three scalar
uppers in \eqref{eq:reference-coefficient-nine:data}.

\subsubsection{The actual finite envelopes}
\label{sec:q117-unified:envelopes}

Use the original variables and receiving weights of
Lemma~\ref{lem:reference-common:receiving}. Every table entry has
\(Z=z_0+Q(B)u\), so its exact mean is \(z_0-\alpha\), where
\(\alpha=-\E[Q(B)u]\). Its mass, variance, and signed covariance
are \eqref{eq:reference-common:mass},
\eqref{eq:reference-common:variance}, and
\eqref{eq:reference-common:covariance}. The fourth-moment bound
\eqref{eq:reference-common:two-endpoints} is used with \(s=m^2\);
in particular its outside scale is \(m^8\).

For the terminal \(qf_b\), write
\[
 \mathsf P_{L,b}(z)=G_L(z)+kz^{L+1}G_{b-1}(-z),\qquad
 k=m^{-1},\quad m=1-\beta_b/2\quad(b\ge3).
\]
For \(b=2\), take \(m=k=1\). All these are finite polynomials.
Write \(V_L=1+\mathsf P_{L,b}(B)u\), and put
\[
 \begin{gathered}
 R_1=KG_{b-1}(K)g,\quad \rho=\E[(Kg)R_1],\quad
 \gamma=k\beta_{b+1}-k^2\|R_1\|^2/2,\\
 a_1=k(d-\beta_{b+2}),\quad B_*=\gamma+a_1,\quad
 a_*=d-k\rho,\quad \lambda=\kappa d,\quad e_4=41d/80.
 \end{gathered}
\]
These are the precise retained coefficients of
\eqref{eq:L58:finiteinputs} and \eqref{eq:L58:propagation}.

For the terminal \(h_2f_b\), use its actual positive square
\(F\le(m+R/2)^2\). Write \(R=\mathcal R(B)u\), where
on a rotation block \(\mathcal R(z)=A(x)-zB_0(x)\), \(x=-z^2\).
With \(D_j=\E B^jR\), put
\[
 \delta=kD_0+k^2\|R\|^2/2,\quad
 \varepsilon=\delta-k\E(KR),\quad \nu_0=1+\varepsilon,
\]
\[
 \mathsf U_L=G_L+kz^L\mathcal R(z),\qquad
 \mathsf Q_L=\mathsf U_L+\varepsilon G_{L-2}+\delta z^{L-2},
 \quad L\ge2.
\]
The last expression equals
\(\mathsf U_L+\delta G_{L-1}-k\E(KR)G_{L-2}\), by the finite
identity \(G_{L-1}=G_{L-2}+z^{L-2}\). Also put
\(\ell=d+k\rho_0\), \(\rho_0=-\E[(Kg)R]\).

The actual square packets have four parameter choices:
\[
\begin{array}{c|c|c}
 b&d\text{ range}&m\text{ and }R\\\hline
 3&(0,1/3]&m=1-d/2,\quad R=K^2g\\
 4&(0,1/4]&m=1-d,\quad R=K^2g+K^3g-dg\\
 b\ge5&(0,1/5]&m=1-9d/8,\quad R=R_{9/4,b}\\
 b\ge5&[1/5,1/4]&m=1-d,\quad R=R_{2,b}.
\end{array}
\]
To specify the last two finite vectors, let \(c=9/4\) or \(2\).
These quotient definitions concern \(d>0\); the degenerate model is
handled directly by leaf integration.
For \(b=5\), take \(C_1=d\), \(C_2=(c-2)d+\delta_2\),
\(C_3=0\). For \(b\ge6\), take
\[
 C_1=\beta_{b-1},\qquad
 \Lambda=\frac{cd-\beta_b-\beta_{b-1}}
 {\beta_{b-2}+(1-d)\beta_{b-3}},\qquad
 C_2=\Lambda\beta_{b-2},\quad C_3=\Lambda\beta_{b-3}.
\]
On the indicated domains these are retained portions of the actual
positive convolution, with \(0\le\Lambda\le1\). Put
\[
 \begin{gathered}
 C_*=C_1+C_2+C_3,\quad D_*=C_2+C_3,\quad E_*=C_3,\\
 R_{c,b}=\sum_{j=3}^bK^j1-C_*g-D_*Kg-E_*K^2g.
 \end{gathered}
\]
The exact identity
\(\beta_b+C_1+C_2+(1-d)C_3=cd\) gives
\(0\le h_2+f_b\le2m+R_{c,b}\). The finite retained-coefficient
proofs are \eqref{eq:L73:retained} for \(c=9/4\) on the low domain
and Section~\ref{loc:S25} for \(c=2\); the latter has its full
retained construction and its stated middle-domain analytic bounds.
The finite norm and coordinate inputs are
\eqref{eq:L77:positivecoordinates},
\eqref{eq:L77:correctedfrequency}, and
\eqref{eq:L77:middleinputs}. Only those actual analytic inputs,
not their genuine-\(f_2\) fourth-moment conclusions, are used here.

\begingroup\small\setlength{\tabcolsep}{3pt}
\begin{longtable}{p{0.17\linewidth}p{0.13\linewidth}p{0.24\linewidth}p{0.28\linewidth}p{0.08\linewidth}}
\caption{Exact envelopes \(Z=z_0+Q(B)u\).
The original comparison is \(0\le T^LF\le m^2Z\) in every row.}
\label{tab:q117-unified:finite-envelopes}\\
\hline
Terminal and range & Scale & \(z_0\) & \(Q(z)\) & Receiver \(t\)\\
\hline\endfirsthead
\hline
Terminal and range & Scale & \(z_0\) & \(Q(z)\) & Receiver \(t\)\\
\hline\endhead
\(qf_b,\ b\ge3,\ L\ge5\)
 & \(1-\beta_b/2\) & \(1-2\lambda\)
 & \(\mathsf P_{L,b}-\lambda\) & \(0\)\\
\(qf_b,\ b\ge3,\ L=4\)
 & \(1-\beta_b/2\) & \(1-\lambda-e_4\)
 & \(\mathsf P_{4,b}-e_4\) & \(0\)\\
\(qf_b,\ b\ge3,\ L=3\)
 & \(1-\beta_b/2\) & \(1-a_*-B_*\)
 & \(\mathsf P_{3,b}-\gamma z-B_*\) & \(t_a\)\\
\(qf_b,\ b\ge3,\ L=2\)
 & \(1-\beta_b/2\) & \(1-\gamma-a_1\)
 & \(\mathsf P_{2,b}-\gamma\) & \(0\)\\
\(qf_b,\ b\ge3,\ L=1\)
 & \(1-\beta_b/2\) & \(1-\gamma\)
 & \(\mathsf P_{1,b}\) & \(0\)\\
\(qf_2,\ L\ge6\)
 & \(1\) & \(1-\tau(4-d)\)
 & \(\mathsf P_{L,2}-\tau(3+2z+z^2)\) & \(0\)\\
\(h_2f_b,\ b\ge3,\ L\ge4\)
 & Square packet & \(\nu_0-2hd\)
 & \(\mathsf Q_L-hd\) & \(0\)\\
\(h_2f_b,\ b\ge3,\ L=2\)
 & Square packet & \(\nu_0\)
 & \(\mathsf Q_2\) & See below\\
\(h_2f_b,\ b\ge3,\ L=3\)
 & Square packet & \(\nu_0-\ell\)
 & \(\mathsf Q_3\) & \(0\)\\
\hline
\end{longtable}
\endgroup

Here \(t_a=0\) for \(a=3\), and \(t_a=\beta_{a-1}\) for
\(a\ge4\). The \(h_2f_b,L=2\) row uses \(t=t_a\) at \(b=3\)
and \(t=0\) at \(b\ge4\). The square packet is the actual
\(b=3\), \(b=4\), low \(b\ge5\), or middle \(b\ge5\) packet,
on its stated continuous domain; changing that analytic envelope leaves \(F\) and the original direction unchanged.

The rows are algebraic forms of the finite positive propagation
proofs \eqref{eq:L58:propagation}, \eqref{eq:L77:propagation}, and
\eqref{eq:L73c:losses}, with their actual finite inputs. In the first two rows the
last two retained losses give respectively
\(Z=V_L-\lambda(1+q)\) and \(Z=V_4-\lambda-e_4q\).
The third row is \(Z=V_3-\gamma Kg-a_*-B_*q\).
The next two rows are \(Z=V_2-\gamma q-a_1\) and
\(Z=q-kKR_1-\gamma\).

For \(qf_2\), the four actual retained path losses have sum
\[
 \bar f_0+\bar f_1+\bar f_2+\bar f_3
 =(4-d)+(3+2B+B^2)u.
\]
This gives its displayed scalar and polynomial, retaining the entire
constant \(4-d\) before multiplying by \(-\tau\).
The \(h_2f_b\) tail subtracts
\(hd(1+q)\), with \(h=9/16-d/2\) at \(b=3\) and
\(h=1/2-6d/5\) for the other low long-connector packets;
its two short rows retain the actual scalar
\(\ell\) at \(L=3\). Thus all means, centered variances, and
receiving covariances follow from the same finite identities, with
their scalar losses and signs retained.

\subsubsection{The finite parameter rows}
\label{sec:q117-unified:row-data}

For \(d>0\), put \(v=\delta_2/d\); then
\(d\le v\le\theta\). In Tables~\ref{tab:q117-unified:short1-data}--
\ref{tab:q117-unified:short2-three-data}, use the normalization
\begin{equation}\label{eq:q117-unified:normalized-data}
 H_r=1-h_rd,\quad U_r=1-a_rd,\quad J_r=dj_r,\quad V_r=dv_r,
 \qquad
 W_r=H_rU_r^4+4J_rU_r^3+M_rV_r.
\end{equation}
The preceding finite receiving identities and the proofs below
give the two actual mean endpoints \(M_rV_r\) and \(W_r\).
Rows marked \(\dagger\) compare the actual receiving prefix after
its own endpoint reduction. The row marked \(\ddagger\) is formed
after the joint covariance--variance reduction in
\eqref{eq:q117-unified:joint-low23}; its displayed \(j_r\) is part
of that complete endpoint polynomial.

The two-step short-one functions are
\begin{equation}\label{eq:q117-unified:two-step-data}
\begin{aligned}
 A_2&=2-3v/2+9v^2/40,& g_2&=A_2-1,\\
 B_2&=2-2v+v^2,& M_2&=1-2v+7v^2/8,\\
 M_{2,f}&=1-2v+2v^2/3,& H_2&=5v/4-17v^2/20,\\
 J_2&=1+v-63v^2/32-855v^4/1024-v^5,\\
 a_0&=1-6v/5+57v^2/80,& e_0&=1-39v/80+57v^3/80,\\
 D_2&=a_0^2+2g_2e_0,&
 \mathcal V_2&=J_2-D_2d+g_2^2d^2.
\end{aligned}
\end{equation}
For the long \((1,2)\) row write
\[
 T_6=(1-v)[2-25v/32+2d(1-v)],\qquad
 \mathcal V_6=1-7v/15-d(1-v)[6-9d-(4+d)v+2v^2].
\]

\begin{table}[htbp]\centering\small
\setlength{\tabcolsep}{4pt}
\caption{Endpoint polynomials for the short-one terminals.}
\label{tab:q117-unified:short1-data}
\begin{tabular}{@{}cccccc@{}}
\toprule
Row&\(h_r\)&\(a_r\)&\(j_r\)&\(v_r\)&\(M_r\)\\
\midrule
S1&\(2\kappa\)&\(3\kappa\)&\(\kappa(1-\kappa)d\)
 &\(1-3\kappa^2d+\kappa^2d^2\)&9\\
S2&\(2\kappa\)&\(27/16\)&\((41/80-\kappa^2)d\)
 &\(1-13d/8+d^2/4\)&9\\
S3L&\(7/4\)&\(13/8\)&\(-1/25+81d/40+57d^2/80\)
 &\(1-3d/2+2d^2-7d^3/5\)&9\\
S3H&\(7/4\)&\(3/2\)&\(-1/25+81d/40+57d^2/80\)
 &\(1-3d/2+2d^2-7d^3/5\)&9\\
S4\(^{\dagger}\)&\(B_2\)&\(A_2\)&\(-M_2+H_2d+3d^2/5\)
 &\(\mathcal V_2\)&9\\
S4F\(^{\dagger}\)&\(B_2\)&\(A_2\)&\(-M_{2,f}+H_2d+3d^2/5\)
 &\(\mathcal V_2\)&9\\
S5&\(2(1-v)\)&\(1-779v/960\)
 &\(-1+11v/5-57v^2/80-\kappa v(1-v)d\)
 &\(1-41v/48-\kappa^2v^2d\)&\(37/4\)\\
S5F&\(2(1-v)\)&1&\(-1+11v/5-v(1-v)d\)
 &\(1-7v/5-v^2d\)&10\\
S6&\(2(1-v)\)&\((1-v)(5-3d)\)&\(-1/10+dT_6\)
 &\(\mathcal V_6\)&9\\
\bottomrule
\end{tabular}
\end{table}

For the short-two long rows put
\(\mathcal W(N,\eta)=N^2-3\eta^2d+\eta^2d^2\),
\(\eta_3=9/16-d/2\), and \(\eta_*=1/2-6d/5\).
\begin{table}[htbp]\centering\small
\caption{Endpoint polynomials for the short-two long connectors.}
\label{tab:q117-unified:short2-tail-data}
\begin{tabular}{@{}cccccc@{}}
\toprule
Row&\(h_r\)&\(a_r\)&\(j_r\)&\(v_r\)&\(M_r\)\\
\midrule
D4F&\(19/16\)&\(27/16-3d/2\)&\(-1/10+4d/5\)
 &\(\mathcal W(1,\eta_3)\)&10\\
D4B&\(19/16\)&\(4/3-18d/5\)&\(3d\)
 &\(\mathcal W(1+d,\eta_*)\)&13\\
D4T&\(19/16\)&\(9/10-18d/5\)&\(3d\)
 &\(\mathcal W(1+7d/5,\eta_*)\)&13\\
\bottomrule
\end{tabular}
\end{table}

Let \(K_*(d)=124/15+38d\), and put
\(C_{23}=1-2v+49v^2/100\), \(A_{23}=1+51v/64\).
The strengthened receiving row D2F covers \(a\ge4\), using the
actual \(a=4\) comparison proved below. D2F3L and D2F3H are the
remaining actual \(a=3\) prefix.
\begin{table}[htbp]\centering\small
\caption{Endpoint polynomials for the short-two two-step connector.}
\label{tab:q117-unified:short2-two-data}
\begin{tabular}{@{}cccccc@{}}
\toprule
Row&\(h_r\)&\(a_r\)&\(j_r\)&\(v_r\)&\(M_r\)\\
\midrule
D2F\(^{\dagger}\)&\(3-3v+v^2\)&\(A_{23}\)&\(-C_{23}+vd\)
 &\(1+v-7v^2/5-d\)&10\\
D2F3L\(^{\ddagger}\)&2&\(A_{23}\)&\(-1+2v+12v^2/11-v^3-d\)
 &\(1+v-7v^2/4-d\)&10\\
D2F3H&2&1&\(-1/20-d\)&\(33/28-d\)&10\\
D2BL\(^{\dagger}\)&\(19/16\)&1&\(-3/16+7d/32\)
 &\(13/10+11d^2/10\)&\(K_*\)\\
D2BM\(^{\dagger}\)&\(19/16\)&1&\(-3/16+7d/32\)
 &\(13/10+11d^2/10\)&13\\
D2TL\(^{\dagger}\)&\(19/16\)&\(7/10\)&\(-1/5+3d/4\)
 &\(45/32+7d/5+36d^2/25\)&\(K_*\)\\
D2TM\(^{\dagger}\)&\(19/16\)&\(2/3\)&\(-1/5+3d/5\)
 &\(45/32+5d/3+16d^2/9\)&14\\
\bottomrule
\end{tabular}
\end{table}

\begin{table}[htbp]\centering\small
\caption{Endpoint polynomials for the short-two three-step connector.}
\label{tab:q117-unified:short2-three-data}
\begin{tabular}{@{}cccccc@{}}
\toprule
Row&\(h_r\)&\(a_r\)&\(j_r\)&\(v_r\)&\(M_r\)\\
\midrule
D3F&\(2(1-v)\)&\(1+25v/64\)&\(-1+3v-v^2-v^3+d/10\)
 &\(1-6v/7-16d/25\)&11\\
D3B&\(19/16\)&\(7/4\)&\(1/10+2d\)
 &\(1+9d/8+3d^2/4\)&11\\
D3TL&\(19/16\)&\(5/4\)&\(1/10+2d\)
 &\(1+5d/2+2d^2\)&12\\
D3TM&\(19/16\)&\(5/4\)&\(1/10+2d\)
 &\(1+8d/3+9d^2/4\)&12\\
\bottomrule
\end{tabular}
\end{table}

\subsubsection{Finite data for the short-one rows}
\label{sec:q117-unified:short-one-data}

Write \(G_a(z)=u_a(x)+zv_a(x)\), where \(z=i\sqrt x\).
The finite alternating coefficients satisfy
\(1-x\le u_a,v_a\le1\) for every \(a\ge3\).
The receiving-symbol lemma above gives
\(\operatorname{Re}(G_a\mathsf P_{L,b})\ge\phi(x)>1/25\)
for every finite \(a,b,L\) in its stated ranges.
For \(a\ge3\), positivity and the actual path convolution give
\begin{equation}\label{eq:q117-unified:receiver-profile}
 0\le f_a\le C_{a,t}\le C_{a,0}
 \le p+\frac65\sqrt d-\kappa d.
\end{equation}
Indeed \(S_a-p=KG_{a-1}(K)g\), and
\(\|G_2(K)\|^2\le45/32<36/25\), whereas
\(\|G_{a-1}(K)\|\le1\) for \(a\ge4\).
Also \(\beta_a\ge\kappa d\). In particular the actual
receiving weights are bounded by \(5/2-q\) on \(d\le1/3\)
and by \(49/20-q\) on \(d\le1/5\).
The increment \(6y/5-\kappa y^2\), \(y=\sqrt d\), is increasing
on the relevant interval. At its two right endpoints use respectively
\(y<26/45\) and \(y<9/20\) to obtain the asserted bounds.

\paragraph{The longer connectors S1--S3.}
For S1, the actual finite propagation
\eqref{eq:L58:propagation} has nonnegative coefficients
\(\gamma,a_1,\varepsilon_j\), with \(\varepsilon_j\ge\tau\)
for \(j\ge4\). Retaining its last two terms gives the first
envelope in Table~\ref{tab:q117-unified:finite-envelopes}.
The finite inputs \eqref{eq:L58:finiteinputs} give
\[
 \|\mathsf P_{L,b}(B)\|\le1,\qquad
 \alpha_L:=-\E[\mathsf P_{L,b}(B)u]\ge\tau,\qquad
 \eta_L:=\langle u,\mathsf P_{L,b}(B)u\rangle\ge\tau,
 \quad L\ge5.
\]
Thus the exact variance after subtracting \(\lambda u\) is at most
\(d-3\kappa^2d^2+\kappa^2d^3\). Its mean is at most
\(1-3\kappa d\). The common covariance identity and
\(\int u_a\,d\nu\le d\) give
\[
 \operatorname{Cov}(C_{a,0},Z)
 =-\int\operatorname{Re}(G_a\mathsf P_{L,b})\,d\nu
   -\beta_{a+1}\alpha_L+\lambda\int u_a\,d\nu
 \le\kappa(1-\kappa)d^2.
\]
These are the S1 data, including mass \(1-2\kappa d\).

For S2 retain the actual coefficient \(\varepsilon_3\ge41d/80\).
The finite fourth-connector calculation following
\eqref{eq:L58:finiteinputs} gives
\[
 \mu\le1-17d/10\le1-27d/16,
 \qquad
 \operatorname{Var}Z\le d-13d^2/8+d^3/4.
\]
In the variance, increasing the coefficient of \(q\) from \(d/2\)
to \(41d/80\) decreases the bound: its derivative is
\(2ed-2\eta_4<0\) for \(e\le41d/80\), since
\(\eta_4\ge\kappa d\). The same covariance calculation gives
\(J=(41/80-\kappa^2)d^2\).

For S3 the finite third-connector inputs of
Section~\ref{proof:Q5-C058} give, with
\(\widehat V=V_3-\gamma Kg\),
\[
 \begin{gathered}
 5d/8\le B_*\le3d/2,\qquad
 1-\E\widehat V\ge(4/5-6d/5)d,\qquad
 \tau\le\alpha_4\le d,\\
 \mu\le1-13d/8\quad(d\le1/4),\qquad
 \mu\le1-3d/2\quad(d\le1/3),\\
 \operatorname{Var}Z\le d-3d^2/2+2d^3-7d^4/5.
 \end{gathered}
\]
Here \(\alpha_4=-\E[\mathsf P_{4,b}(B)u]\).
Use \(t=t_a\) from Table~\ref{tab:q117-unified:finite-envelopes}.
Its actual mass is \(1-2d\) for \(a=3\), and at most
\(1-3\tau\) for \(a\ge4\), hence at most \(1-7d/4\).
Let
\(M_a=\int\operatorname{Re}(G_a\mathsf P_{3,b})\,d\nu\),
\(J_a=\int xv_a\,d\nu\), and \(D_a=\int u_a\,d\nu\).
The covariance, with every receiving correction retained, is
\begin{equation}\label{eq:q117-unified:third-shortone-covariance}
 \operatorname{Cov}(C_{a,t_a},Z)
 =-M_a-\gamma J_a
  -\beta_{a+1}(1-\E\widehat V)
  +B_*D_a+t_a\alpha_4-t_aB_*d.
\end{equation}
Use \(M_a\ge d/25\), \(J_a\ge0\), \(D_a,t_a\le d\),
and the preceding finite bounds. Dropping the nonpositive terms gives
\[
 \operatorname{Cov}(C_{a,t_a},Z)
 \le-d/25+(81/40)d^2+(57/80)d^3.
\]
This proves S3L and S3H on their respective continuous domains.

\paragraph{The two-step terminal data S4 and S4F.}
Put \(V_R=\|R_1\|^2\), \(\rho=\langle Kg,R_1\rangle\),
\(D_3=\E B^3R_1\), \(t_0=\beta_{b+1}\),
and \(e=d-\beta_{b+2}\). The actual finite identities
\eqref{eq:L58:two}--\eqref{eq:L58:twomoments} are
\begin{equation}\label{eq:q117-unified:two-step-identities}
\begin{aligned}
 \alpha&=d-k\rho,& \eta&=d-kD_3,\\
 \gamma&=kt_0-k^2V_R/2,& a_1&=ke,\\
 \mu&=1-\alpha-\gamma-a_1,&
 \operatorname{Var}Z
  &=\|V_2-1\|^2-\alpha^2+\gamma^2d-2\gamma\eta.
\end{aligned}
\end{equation}
In particular \(\rho=d-\beta_{b+3}\), and the finite alternating
sums give
\[
 t_0+e-\rho=
 \begin{cases}\beta_{b+1},&b\text{ odd},\\
                \beta_{b+3},&b\text{ even}.
 \end{cases}
\]
Thus \(1-\mu=d+k(t_0+e-\rho)-k^2V_R/2\).
The finite geometric loss \eqref{eq:Q81:finite-geometric-loss}
implies the weaker \(\|G_N(i\sqrt x)\|^2\le1-9x/20\)
for every \(N\ge3\). For \(b\ge4\),
\(V_R\le\delta_2-9\delta_3/20\), and
\(k\tau-k^2V_R/2\) increases on \(1\le k\le6/5\),
since its derivative is at least \(17d/160\).
Jensen's moment inequality therefore gives
\(\mu\le1-A_2d\) and \(\gamma\ge g_2d\).
For \(b=3\), use instead \(t_0+e-\rho=d\) and
\(V_R=\delta_2+\delta_3\). The same derivative is positive,
and \(1-\mu\ge(2-45v/64)d\),
\(\gamma\ge(1-45v/64)d\), which imply these same two bounds.

The finite norm proof following \eqref{eq:L58:twomoments} gives
\(\|V_2-1\|^2\le dJ_2(v)\), with the exact \(J_2\) in
\eqref{eq:q117-unified:two-step-data}. For \(b\ge5\), the two
finite alternating coefficients of \(G_{b-1}\) are at most
\(1-\kappa x\); hence
\(\rho\le\delta_2-\kappa\delta_3\) and
\(D_3\le\delta_3-\kappa\delta_4\).
Using \(\delta_3\ge v^2d\), \(\delta_3\le\theta vd\),
and \(\delta_4\ge v^3d\) gives
\(\alpha\ge a_0d\) and \(\eta\ge e_0d\).
In \eqref{eq:q117-unified:two-step-identities}, the larger endpoint
of the convex expression in \(\gamma\in[g_2d,6d/5]\) is
\(\gamma=g_2d\), because
\([(6/5)+g_2]d-2e_0<0\). Consequently
\begin{equation}\label{eq:q117-unified:shortone-common-variance}
 \operatorname{Var}Z\le d\mathcal V_2(v,d),\qquad b\ge3.
\end{equation}
The two finite prefixes in this last assertion
use the following comparisons. At \(b=4\), put
\(a_4=1-6v/5+6v^2/5\), \(e_*=1-39v/80\). Then
\[
 a_4^2-a_0^2\ge1599v^2/3200,
 \qquad2g_2(e_0-e_*)\le57v^2/224.
\]
These follow from \(a_4-a_0=39v^2/80\),
\(a_4+a_0\ge41/40\), \(g_2\le1-7v/5\), and
\(v(1-7v/5)\le5/28\). At \(b=3\), put
\(a_*=1-6v/5\), \(g_*=1-45v/64\). The actual comparison is
\(dJ_2-(a_*^2+2g_*e_*)d^2+g_*^2d^3\).
Here \(g_*-g_2\ge7v/10\), \(g_*+g_2\le2-2v\), and
\(2e_*-(g_*+g_2)d>6/5\); the difference from
\eqref{eq:q117-unified:shortone-common-variance} is paid by
\((21/25)v-(1881/1120)v^2>0\).

\paragraph{The two-step receiving symbols and endpoints.}
For \(a\ge4\), every finite terminal satisfies
\[
 \operatorname{Re}(G_a\mathsf P_{2,b})\ge1-2x+2x^2/3;
 \quad b\ge5\ \Longrightarrow\
 \operatorname{Re}(G_a\mathsf P_{2,b})\ge1-2x+7x^2/8.
\]
Here is the finite-index verification. Write
\(G_{b-1}(-z)=\alpha_0-z\beta_0\) and
\(\mathsf P_{2,b}=c+ze\), with
\(c=1-kx^2\beta_0\), \(e=1-kx\alpha_0\).
For \(b\ge5\), \(\alpha_0-\beta_0\ge-x^3\), so
\(c-xe>0\). At \(a=4,5\), use
\(G_4=(1-x)(1+z)\) and \(G_5=G_4+x^2\), together with
\(kx(1-x)\le3/10\). For \(a\ge6\), the finite coefficients
satisfy \(v_a\ge(1-x)(1+x^2)\), \(u_a-v_a\ge-x^3\).
The sufficient excess divided by \(x^2\) is at least
\[
 9/8-3x+7x^2/10-3x^4/10
 \ge142669/10485760>0.
\]
The polynomial decreases on \([0,\theta]\).
For arbitrary \(b\ge3\), use instead
\(\alpha_0-\beta_0\ge-x\),
\(\alpha_0\ge(1-x)\beta_0\), and \(\beta_0\ge1-x\).
At \(a=4\) the payment is \(kx(1-x)\le3/10<1/3\).
For \(a\ge6\) it reduces to
\[
 4/3-3x+x^2-k(1-x)x+k(1-x)x^4>0.
\]
Its minimum in \(k\) occurs at \(6/5\); the resulting polynomial
decreases, and at \(\theta\) its first three combined terms are
\(-151/15360\), while its final term exceeds \(1/60\).
The remaining receiver \(a=5\) has the stronger bound
\(\operatorname{Re}(G_5\mathsf P_{2,b})\ge1-2x+3x^2/2\)
for every \(b\), from \(kx<1/2\).

Let \(M_a=\int\operatorname{Re}(G_a\mathsf P_{2,b})\,d\nu\),
\(D_a=\int u_a\,d\nu\le d-\kappa\delta_2\). The exact identity is
\[
 \operatorname{Cov}(C_{a,0},Z)=-M_a-\beta_{a+1}\alpha+\gamma D_a.
\]
Finite alternating coefficients also give
\(V_R\ge c_G\delta_2\), \(c_G=16245/32768\).
Keeping the same \(k\), the last two terms divided by \(d^2\)
are at most
\[
 (k-1)+v[1+k(1-\kappa)]-kv^2
 -\tfrac12k^2(V_R/d)(1-\kappa v)
 \le H_2(v)+3d/5.
\]
Here \(k-1\le3d/5\),
\(1+(6/5)\theta-c_G/2<5/4\), and
\(1-c_G\kappa/2>17/20\). For \(a=3\), its actual mass is
\(1-2d\), and \(-d\alpha+\gamma\tau\le(k-1)d^2\), so the
stronger covariance is \(-(1-2v)d+3d^3/5\).

For \(b\ge5\), the mass \(H=1-B_2d\) and covariance in row S4
therefore apply directly at \(a=4\) and \(a\ge6\).
The finite mass losses retain \(2\tau+\delta_3\): this is exact
at \(a=6\); at \(a=4\) use \(\delta_2\ge\delta_3\); and at
\(a\ge7\) each sum retains \(\tau+\kappa\delta_3\), with
\(2\kappa>1\). Compare the remaining receivers only after applying
their own two mean endpoints. Their zero endpoints are unchanged.
At \(a=5\), the nonzero endpoint difference from S4 is at most
\(v^2dU^3(U-5/2)\le0\). At \(a=3\), divided by \(dU^3\), it is
\[
 v[-2+9v/2]+dv[-1-8v/5+39v^2/20-9v^3/40]\le0,
\]
since \(v\le\theta<4/9\) and \((39/20)\theta<8/5\).
This proves the complete S4 endpoints for every receiving length.
For \(b=3,4\), replace \(M_2\) by \(M_{2,f}\) as in S4F.
The same receiving comparisons remain valid. Its nonzero endpoint
increases by \((5/6)v^2dU^3\), while the actual inverse scale obeys
\begin{equation}\label{eq:q117-unified:shortone-prefix-scale}
 (1-d/2)^{-8}-[1-(1-v)d/2]^{-8}
 \ge4vd\ge(5/6)v^2dU^3.
\end{equation}
The first inequality follows termwise from the positive inverse
binomial series. Thus the S4 certificate also pays S4F.

\paragraph{The unit connector S5 and S5F.}
For \(b\ge4\), the actual \(Z=q-kKR_1-\gamma\) satisfies
\[
 \mu=1-k(t_0+e)+k^2V_R/2,\qquad
 \operatorname{Var}Z=d-2k\rho+k^2\|KR_1\|^2-k^2e^2.
\]
By \eqref{eq:Q81:finite-geometric-loss},
\(V_R\le\delta_2-7\delta_3/15\) and
\(\|KR_1\|^2\le\delta_3-7\delta_4/15<\delta_2/3\).
The finite alternating identities give
\(\rho,e\ge\kappa\delta_2\) and \(t_0+e\ge d-\delta_3\).
The last sum equals \(d\) for even \(b\); its first negative
odd correction is \(-\delta_3\), and later ones are smaller.
Both derivatives in \(k\) have the favorable sign:
\(t_0+e-kV_R>0\) and \(k\|KR_1\|^2<\rho\).
At \(k=1\) these yield exactly the S5 mean and variance.
For its receiving symbol write \(G_{b-1}(-z)=\alpha_0-z\beta_0\).
Every finite \(b\ge4\) satisfies
\((1-x)\alpha_0+x\beta_0\le1-\kappa x\).
Since the two coefficients of \(\mathsf P_{1,b}\) are nonnegative,
\[
 \operatorname{Re}(G_a\mathsf P_{1,b})
 \ge1-11x/5+57x^2/80.
\]
The common covariance identity, with
\(\alpha_1=ke\ge\kappa\delta_2\), gives the S5 value of \(J\).
For \(b=3\), the exact data are
\(R_1=Kg+K^2g\), \(t_0=d\), \(e=\rho=\delta_2\),
\(V_R=\delta_2+\delta_3\), and
\(\|KR_1\|^2=\delta_3+\delta_4\).
The same derivative argument gives
\(\mu\le1-d\) and
\(\operatorname{Var}Z\le d(1-7v/5)-v^2d^2\).
Here \(2-\theta-\theta^2>7/5\), and the finite receiving symbol is
\(\operatorname{Re}(G_a\mathsf P_{1,3})\ge1-11x/5\).
Together with \(\alpha_1=k\delta_2\), these prove S5F.

\paragraph{The four-loss tail S6.}
The exact identity
\(qf_2=(1+Kg/2)^2-(g+Kg/2)^2\), positivity of \(T\),
and \(T((Kg)^2)\le2\delta_2\) give the finite propagation
\[
 0\le T^L(qf_2)\le V_L-\gamma\bar f_{L-1}
       -\sum_{j=3}^L\omega_j\bar f_{L-j},\qquad
 \gamma=d-\delta_2/2,\quad
 \omega_j=\begin{cases}\beta_j,&j\text{ even},\\
                         \beta_{j+2},&j\text{ odd}.
             \end{cases}
\]
This is the finite Duhamel expansion of the actual \(Kg\) term;
every coefficient \(\omega_j\ge\tau\). For \(L\ge6\), the last
four indices are at least three and give the S6 mask. Its unmodified
mean shift is \(\alpha_L=\beta_{L+3}\) for even \(L\),
\(\beta_{L+1}\) for odd \(L\), hence in \([\tau,d]\).
The four retained path means sum to \(4-3d\), proving
\(\mu\le1-\tau(5-3d)\).

For every finite \(n\ge4\), the receiver-independent finite
calculation \eqref{eq:L33:finiteP} gives
\( |\mathsf P_{n,2}|^2\le1-7x/15\).
For \(L\ge6\), \eqref{eq:L33:RI} gives both finite coefficients
in \([\rho_*(x),1]\), where
\(\rho_*(x)=1-x+x^2-x^3\).
Those proofs include the four residue classes and their actual
finite last exponents.
The actual receiving symbol consequently obeys
\(\operatorname{Re}(G_a\mathsf P_{L,2})\ge7/60\).
Indeed its finite lower value is at least
\(B_0-D_0x^4\), with
\(B_0=(1-2x)/(1+x)\), \(D_0=(1-x+x^2)/(1+x)\).
This decreases, and at \(\theta\) exceeds
\(2/15-1/60=7/60\), because
\(259\cdot28561<8388608\).

In the S6 mask the real and imaginary coefficient subtractions
are \(\tau(3-x)\) and \(2\tau\). The componentwise proof of
\eqref{eq:L33:variance} permits each of its four actual retained
coefficients to lie in \([\tau,d]\). Setting all four equal to
\(\tau\) gives precisely this S6 polynomial. Its same-measure
spectral integration, using \(\delta_3\ge v^2d\) and projection
onto mean-zero functions, gives the stated S6 variance.
The actual receiving covariance is computed separately below.
For the covariance let \(D_a=\int u_a\,d\nu\),
\(J_a=\int xv_a\,d\nu\), \(E_a=\int xu_a\,d\nu\).
The actual centered mean shift is \(\alpha_L-2\tau d\), so
\[
 \operatorname{Cov}(C_{a,0},Z)
 =-M_a+\tau(3D_a-2J_a-E_a)
       -\beta_{a+1}(\alpha_L-2\tau d).
\]
On \(d\le5/24\), this last mean shift is positive. Use
\(D_a\le d\), \(J_a,E_a\ge\kappa\delta_2\),
\(\beta_{a+1},\alpha_L\ge\tau\), and \(M_a\ge7d/60\).
Since \(3\kappa-1=25/32\), the resulting bound is exactly
the S6 value of \(J\), after weakening \(-7d/60\) to \(-d/10\).

\subsubsection{Finite data for the short-two rows}
\label{sec:q117-unified:short-two-data}

The four actual square packets in
Table~\ref{tab:q117-unified:finite-envelopes} have the following
finite analytic bounds. They are the propagation, coordinate, and
moment bounds of \eqref{eq:L77:propagation},
\eqref{eq:L77:twomoments}--\eqref{eq:L77:threemoments},
\eqref{eq:L77:positivecoordinates},
\eqref{eq:L77:correctedfrequency}, and
\eqref{eq:L77:middleinputs}, together with the long-connector
finite loss calculation \eqref{eq:L73c:losses}.
Here \(N\) bounds both \(\|\mathsf Q_L\|\) and
\(\|\mathsf Q_{L-1}\|\) when \(L\ge4\), and
\(h d\) bounds the retained losses and the three actual mean
shifts at those connector lengths.
\begin{equation}\label{eq:q117-unified:square-packet-data}
\begin{array}{c|c|c|c|c|c}
 b& d\text{ range}&\varepsilon&\delta&h&N\\\hline
 3&(0,1/3]&-\eta&k^2\delta_3/2&9/16-d/2&1\\
 4&(0,1/4]&|\varepsilon|\le d/6&\le25d/24&1/2-6d/5&1+d\\
 b\ge5&(0,1/5]&|\varepsilon|\le3d/5&\le6d/5&1/2-6d/5&1+7d/5\\
 b\ge5&[1/5,1/4]&|\varepsilon|\le2d/3&\le4d/3&\text{not used}&\text{not used}
\end{array}
\end{equation}
At \(b=3\), \(\eta=k\delta_2-k^2\delta_3/2\),
\(0\le\eta\le39d/80\), and \(0\le\delta\le\eta/3\).
For \(b=4\), \(A=d+x\), \(B_0=x\).
For \(b\ge5\), both coordinates \(A,B_0\) are nonnegative.
On the low domain \(\sqrt{A^2+xB_0^2}\le3/8+9d/5\), and
\(k(3/8+9d/5)\le1\); the \(b=4\) packet also has
\(k\sqrt{A^2+xB_0^2}\le1\).
These are bounds at fixed original moments: in their finite proof
one first maximizes the actual functions of \(x\), and only then
uses \(\delta_2\le\theta d\). No substitution \(\delta_2=xd\)
is made in a frequency polynomial.

For all three low long-connector packets, the positive propagation
and the actual mean shifts give
\begin{equation}\label{eq:q117-unified:long-square-moments}
 \mu\le\nu_0-3hd,\qquad
 \operatorname{Var}Z\le N^2d-3h^2d^2+h^2d^3,
 \qquad
 Z\le(\nu_0-hd)q-hd+N\sqrt d.
\end{equation}
The negative square comes from the actual mean of
\(\mathsf Q_L(B)u\), as in
\eqref{eq:reference-common:variance}.

\paragraph{Long connectors D4F, D4B, and D4T.}
First let \(b\ge4\) in its low packet. The finite polynomial
\(\mathsf U_L=G_L+kz^L(A-zB_0)\) satisfies
\[
 \operatorname{Re}(G_a\mathsf U_L)\ge1/50,
 \qquad a\ge3,\quad L\ge4.
\]
At \(L=4\) the perturbation contributes
\(kx^2(u_a A+xv_aB_0)\ge0\), and the unperturbed product is
at least \((1-x)(1-2x)\ge57/512\).
At \(L=5\) its unperturbed value is at least
\(B_0^*+A_0^*x^3\ge2/15\), with
\(B_0^*=(1-2x)/(1+x)\), \(A_0^*=(2-x)/(1+x)\).
The perturbation costs at most \(x^{5/2}\le1/9\), leaving
\(1/45>1/50\). At \(L=6\) the unperturbed product exceeds
\(B_0^*\); at \(L=7\) it is at least
\(2/15-(17/15)\theta^4>11/108\).
For \(L\ge8\), the finite even and odd formulas give, respectively,
\((2/15)(1-\theta^4)>7/54\) and
\(2/15-2\theta^6>67/540\), both exceeding \(3/25\).
The perturbation at \(L\ge6\) costs at most \(\theta^3<7/100\).
These cover every finite connector length.

In \(\mathsf Q_L=\mathsf U_L+\varepsilon G_{L-2}
 +\delta z^{L-2}\),
\(|\operatorname{Re}(G_aG_{L-2})|\le1\); at \(L=4\) this
follows directly from \(0\le u_a-xv_a\le1\), and later from
the two finite norm bounds. The last correction is bounded below
by \(-\delta\theta\). Thus the additional receiving cost is at
most \(151d/256\) for \(b=4\) and \(87d/80\) for \(b\ge5\).
Using the exact centered mean \(\alpha_L\ge hd\), the covariance is
\begin{equation}\label{eq:q117-unified:long-square-covariance}
 \operatorname{Cov}(C_{a,0},Z)
 =-\int\operatorname{Re}(G_a\mathsf Q_L)\,d\nu
       -\beta_{a+1}\alpha_L+hd\int u_a\,d\nu
 \le3d^2.
\end{equation}

For \(b=3\), \(\mathsf U_L=G_L-kz^{L+2}\) has the stronger
finite receiving lower bound \(1/10\). The \(L=4\) correction
is favorable. At \(L=5\), its cost \(kx^4\) is paid by
\(A_0^*x^3\), since \(A_0^*\ge17/15>kx\).
At \(L=6\) the lower value is
\(B_0^*+x^3(B_0^*-kx)\ge2/15-(17/48)(7/100)>1/10\).
At \(L=7\) the correction is favorable and the unperturbed
value exceeds \(11/108\). For \(L\ge8\) the cost is at most
\((6/5)\theta^5<1/50\), paid by the unperturbed \(3/25\).
The remaining corrections cost at most
\(\eta(1+\theta/3)\le4251d/7680\).
Retaining \(-\beta_{a+1}\alpha_L\), rather than discarding it,
now gives
\[
 \operatorname{Cov}(C_{a,0},Z)
 \le-d/10+[4251/7680+\theta h]d^2
 \le-d/10+4d^2/5,
\]
since \(4251/7680+\theta(9/16)=1001/1280<4/5\).
Together with \eqref{eq:q117-unified:long-square-moments}, these
are D4F, D4B, and D4T. In every case the actual receiving mass
is at most \(1-2\tau\le1-19d/16\).

\paragraph{The finite terminal \(b=3,L=2\).}
For this packet the actual polynomial and scalar are
\[
 \mathsf Q_2=1+z+\delta-kz^4,\qquad
 Z=1-\eta+\mathsf Q_2(B)u,\qquad \mu=1-d-\eta.
\]
Since \(k\delta_2-k^2\delta_3/2\) increases with \(k\),
\(\eta\ge\delta_2-\delta_3/2\ge51\delta_2/64\).
The exact variance of this same message is
\[
 (1+\delta)^2d+\delta_2-2k(1+\delta)\delta_3
       +k^2\delta_5-d^2.
\]
It is at most \(d(1+v-7w/4-d)\) on \(d\le1/6\), and at
most \(d(1+v-7v^2/5-d)\) on \(d\le1/3\), where
\(w=\delta_3/d\). To see the first bound, combine the exact
\(\delta_3\) terms before estimating them. After multiplication
by \(m^2\), the sufficient difference is
\(1/4-\theta^2-d/4-7d^2/16\), decreasing and equal to
\(287/9216>0\) at \(1/6\). The full-range coefficient is
\(7/5\), as in the finite calculation for \eqref{eq:L77:23}.

For \(a\ge4\), use \(t=\beta_{a-1}\) and put
\(M_a=\int\operatorname{Re}\{G_a(1+z-z^4)\}\,d\nu\).
The exact receiving cancellation is
\begin{equation}\label{eq:q117-unified:retained-receiving-cancellation}
\begin{aligned}
 \operatorname{Cov}(C_{a,\beta_{a-1}},Z)
 ={}&-M_a+d(\beta_{a-1}-\beta_{a+1})\\
 &+\delta\left(\beta_{a-1}d-\int u_a\,d\nu\right)
 +(k-1)\int u_ax^2\,d\nu-k\beta_{a-1}\delta_3.
\end{aligned}
\end{equation}
The \(\delta\) term is nonpositive, since
\(u_a\ge\kappa>d\ge\beta_{a-1}\). The last two terms have
nonpositive sum: \(k-1=kd/2\), \(u_a\le1\), and
\(\beta_{a-1}\ge\tau>d/2\). The beta difference is one signed
\(\delta_j\), \(j\ge2\). Hence the covariance is at most
\(-M_a+vd^2\).
For \(a=4\), the exact frequency polynomial is
\(1-2x+x^3\), and the actual receiving mass is
\(1-(3-v)d\). For every finite \(a\ge5\),
\[
 M_a\ge(1-2v+49v^2/100)d,\qquad
 \E C_{a,\beta_{a-1}}\le1-(3-3v+v^2)d.
\]
The finite frequency polynomials at \(a=5,6,7\) are respectively
\(1-2x+x^2+x^3-x^4\),
\(1-2x+x^2-x^4\), and
\(1-2x+x^2-x^3-x^4+x^5\).
In the last, the coefficient after \(x^2\) is
\((1-x)(1-x^2)\ge16245/32768>49/100\).
For \(a\ge8\), both coefficients are at least
\((1-x)(1+x^2)\), their difference is at least \(-x^5\),
and the product is at least
\(1-2x+x^2-x^3+x^7\ge1-2x+\kappa x^2\).
Integrate and use \(\delta_3\ge v^2d\).
The three actual beta losses at \(a=5,6,7\) retain
\(3\tau+\delta_3\); at \(a\ge8\), each retains
\(\tau+\delta_3-\delta_4\), and \(3\kappa>1\).
This proves the mass bound as well.

The D2F endpoints therefore apply directly at \(a\ge5\).
At \(a=4\), after taking its own two mean endpoints, the nonzero
endpoint difference from D2F divided by \(dU^3\) is
\[
 v[-(2-v)U+(49/25)v-4v^2]\le0.
\]
Indeed \(U\ge3433/6144>5/9\), and the bracket is at most
\(-(2-\theta)(5/9)+(49/25)\theta=-107/1200<0\).
The zero endpoint is the same. This extends D2F to all \(a\ge4\).

For \(a=3\), the actual mass is \(1-2d\), and the covariance is
\[
 -d+2\delta_2+k\delta_3-k\delta_4
       -\delta(d-\delta_2)-d^2.
\]
On \(d\le1/6\), \(k\le12/11\). With
\(w=\delta_3/d\), \(w_4=\delta_4/d\), the actual joint bounds are
\[
 J\le d[-1+2v+(12/11)w-w_4]-d^2,\qquad
 V\le d[1+v-7w/4-d].
\]
They must be combined before reducing the moments. In the complete
nonzero endpoint, the coefficient of \(w\), divided by \(d\), is
\((48/11)U^3-35/2\le-289/22<0\); the coefficient of \(w_4\)
is \(-4U^3\le0\). Consequently \(w\ge v^2\) and
\(w_4\ge v^3\) give
\begin{equation}\label{eq:q117-unified:joint-low23}
 HU^4+4JU^3+10V\le W_{\mathrm{D2F3L}}(d,v).
\end{equation}
At the zero endpoint, \(10V\) separately decreases with \(w\).
Thus D2F3L is a joint endpoint row, not an independent covariance
upper obtained by substituting \(w\ge v^2\) into \(J\).
On \(d\ge1/6\), use \(w^2\le vw_4\), \(w\le\theta v\).
The quadratic \(-1+2v+(6/5)w-w^2/v\) increases in
\(w\in[0,\theta v]\); its resulting upper increases with \(v\).
Thus \(-1+2v+(6/5)w-w_4\le-9257/163840<-1/20\),
which gives D2F3H with \(\mu\le1-d\) and
\(V\le d(33/28-d)\).

\paragraph{The two-step finite symbols for \(b\ge4\).}
The actual centered mean shift is
\(\alpha_2=d+k\rho_0\), with
\(\rho_0=\int xB_0\,d\nu\ge0\).
For \(b=4\), remove the nonlinear retained constants and set
\(k=1\) to obtain the baseline real coefficients
\(1-x^2\), \(1+x^2\) of \(\mathsf U_2\).
Its receiving lower bound is \(1-2x-x^2>1/50\) at \(a=3\).
For \(a=4\) it is \(1-2x+x^4\ge3/16\), and \(a=5\) is larger.
For \(a\ge6\), the finite lower polynomial is
\(1-2x+x^2-3x^3+x^4+x^5+x^6\).
It decreases on \([0,\theta]\), and its excess over \(3/16\)
at \(\theta\) is \(\theta^2(42545/1048576)>0\).
The baseline vector norm is at most \(1/2\). Since
\(k-1\le4d/3\), the actual perturbation of \(\mathsf U_2\)
has norm at most \(13d/16\). The extra real shift \(\delta\)
has a favorable receiving sign. The common covariance formula
therefore gives
\[
 J_3\le-d/50-3d^2/16,\qquad
 J_a\le-3d/16+7d^2/32\quad(a\ge4).
\]

For \(b\ge5\), write the actual finite decomposition
\(A=A_0+C_*-E_*x\), \(B_0=B_{00}+D_*\).
At \(b=5\), \(A_0=x-x^2\), \(B_{00}=x\).
At \(b\ge6\), \(A_0=xa_f\), \(B_{00}=x\bar a_f\), with
the two finite alternating coefficients
\(a_f,\bar a_f\in[1-x,1-x+x^2]\).
The baseline vector norm is at most \(3/8\): at \(b=5\)
its square is \(x^2-x^3+x^4\), and later it is at most
\((x-x^2+x^3)^2(1+x)<(3/8)^2\).
The receiving baseline obeys
\[
 \operatorname{Re}(G_3\mathsf U_2^0)
 \ge1-2x-x^2+\kappa x^3>3/50,\qquad
 \operatorname{Re}(G_a\mathsf U_2^0)\ge1/5\quad(a\ge4).
\]
For \(a=4,5\), use \(a_f+x\bar a_f\le1+x^3\), so the
lower bound is \((1-x)(1-x-x^2-x^5)>1/5\).
For \(a\ge6\), the finite coefficients give
\(1-2x+x^2-2x^3+x^6-x^7+x^8\); its derivative is negative,
and its endpoint exceeds
\(7/20-7/50-1/500=26/125>1/5\).
The \(a=3\) polynomial decreases and at \(\theta\) is
\(65295/1048576>3/50\).

On the low domain the true retained coefficients give
\(0\le C_*-E_*x\le15d/8\), \(0\le D_*\le101d/80\).
Their perturbing vector norm is at most \(21d/10\), and
\(k-1\le45d/31\), \(k\le40/31\) give a total
\(\mathsf U_2\) perturbation \((10491/7936)d\).
On the middle domain the corresponding constants are
\(C_*\le53d/32\), \(D_*\le17d/16\), vector norm
\(9d/5\), and total perturbation \((377/320)d\).
Including the actual mean product yields
\[
 \begin{array}{c|cc}
 &a=3&a\ge4\\\hline
 d\le1/5&J\le-3d/50+d^2/3&J\le-d/5+3d^2/4\\
 1/5\le d\le1/4&J\le-3d/50+d^2/5&J\le-d/5+3d^2/5.
 \end{array}
\]
The unrounded quadratic coefficients are \(2555/7936,5779/7936\)
and \(57/320,187/320\), respectively.
The exact two-step mean and variance data from
\eqref{eq:L77:twomoments}, \eqref{eq:L77:shiftedPtwo}, and
\eqref{eq:L77:middleinputs} are those in D2BL, D2BM, D2TL, D2TM.
For each \(a=3\) receiver take its own two mean endpoints first.
The nonzero endpoint differences from the formal \(a\ge4\) rows are
\[
 \begin{array}{c|c}
 b=4&-dU^3(57/400+13d/16)\\
 b\ge5,\ d\le1/5&-dU^3(101/400+527d/480)\\
 b\ge5,\ 1/5\le d\le1/4&-dU^3(101/400+127d/120).
 \end{array}
\]
They are nonpositive, and the zero endpoints agree. This proves
the all-receiver endpoint interpretation of those four rows.

\paragraph{Three-step connectors D3B, D3TL, and D3TM.}
The scalar \(\ell=d+k\rho_0\) is retained in the mean.
It cancels from the covariance, which is exactly
\[
 J=-\int\operatorname{Re}(G_a\mathsf U_3)\,d\nu
       -\varepsilon\int u_a\,d\nu
       +\delta\int xv_a\,d\nu-\beta_{a+1}\alpha_3.
\]
For the baseline packets just used, the two coefficients are
\(1-x-x^3\bar a_f\) and \(1-x^2a_f\).
Use \(a_f=\bar a_f=1\) at \(b=4\), and
\(a_f=1-x,\bar a_f=1\) at \(b=5\).
In every finite case \(a_f-(1-x)\bar a_f\ge0\), so
\[
 \operatorname{Re}(G_a\mathsf U_3^0)
 \ge1-3x+x^2+x^3[a_f-(1-x)\bar a_f]
 \ge-55/1024>-3/50.
\]
The three-step perturbation is at most \(\sqrt\theta<13/20\)
times the two-step perturbation. Thus it is below \(9d/10\)
on the low \(b\ge5\) domain, \(4d/5\) on its middle domain,
and \(169d/320\) at \(b=4\).
The actual finite mean shifts satisfy \(\alpha_3\ge\kappa d\).
The resulting covariance fees, divided by \(d^2\), are bounded by
\[
 9/10+3/5+(6/5)\theta-\kappa^2,\quad
 4/5+2/3+(4/3)\theta-\kappa^2,\quad
 169/320+1/6+(25/24)\theta-\kappa^2,
\]
respectively; every one is below two. Hence \(J\le d/10+2d^2\).
The actual means and centered variances of
\eqref{eq:L77:threemoments} and \eqref{eq:L77:middleinputs}
give D3B, D3TL, and D3TM.

\paragraph{The finite three-step terminal D3F.}
For \(b=3\), the exact polynomial is
\[
 \mathsf Q_3=1+z+z^2-\eta+\delta z-kz^5,\qquad
 Z=1-\eta-d+\mathsf Q_3(B)u.
\]
Its centered mean shift is
\(\alpha_3=d-k\delta_3+\delta d\ge2053d/2560>4d/5\), and
\[
 \mu=1-2d-k\delta_2+k\delta_3+(1-d)\delta
 \le1-2d-25\delta_2/64\le1-(1+25v/64)d.
\]
Indeed \(k(1-d)\le1\), and
\(k[\delta_2-(3/2)\delta_3]\ge25\delta_2/64\).
Before the \(\eta,\delta\) shifts, the squared norm is at most
\(1-x+x^2-2x^3+x^5\le1-6x/7\):
\[
 1/7-x+2x^2-x^4\ge1/7-x+(11/6)x^2\ge1/154>0.
\]
The actual shifts change the norm by at most
\[
 -2\eta(1-x)+\eta^2+2\delta x(1-kx^2)+\delta^2x
 \le-11\eta/12+\eta^2(1+\theta/9)\le0,
\]
using \(\delta\le\eta/3\), \(\eta\le13/80\).
Subtracting the true mean square gives
\(V\le d-(6/7)\delta_2-(16/25)d^2\), as in D3F.
The baseline receiving product satisfies
\[
 \operatorname{Re}\{G_a(1+z+z^2-z^5)\}
 \ge1-3x+x^2+x^3.
\]
Its two coordinates are nonnegative; increasing \(k\) only
increases that product. The \(\eta,\delta\) fees are at most
\((109/80)vd^2\), and the negative mean product is at least
\((4\kappa/5)d^2\). Their difference is at most
\((201/2560)d^2<d^2/10\).
Integrating the baseline with \(\delta_3\ge v^2d\) and
\(\delta_4\ge v^3d\) gives exactly the D3F covariance.

\subsubsection{Pointwise fourth coefficients}
\label{sec:q117-unified:point-coefficients}

Write \(y=\sqrt d\), and put
\[
 \mathfrak C=2-q+6y/5-\kappa y^2,\qquad
 \mathfrak C_\infty=5/2-q,\qquad
 \mathfrak C_\ell=49/20-q.
\]
The receiving profiles \eqref{eq:q117-unified:receiver-profile}
give the indicated upper for every actual \(C_{a,t}\).
All three profiles are nonnegative on the rectangles used below.
We prove the coefficient bounds in the row tables by bounding
\(Z+\mu\) and its actual mean \(\mu\), then certifying the
resulting polynomial in \((q,y)\).

\paragraph{The two coefficient-nine entrances.}
S1, S2, S3L, S3H, and S6 satisfy the hypotheses (Q117.4.1).
For S1 and S6 their actual point upper is at most
\((1-\kappa d)q-\kappa d+y\); the difference from the
(Q117.4.1) upper is \(y/12+(3-q)d/32\ge0\).
For S2 the difference is at least \(y/12-21d/160\ge0\),
since \(y\le3/5\). S3 has exactly the required finite point
upper from its third-connector propagation.
Their actual means are at most \(1-3d/2\).
Thus \eqref{eq:q117-unified:weak-nine} gives their coefficient nine.

For clarity retain the stronger two-step entrance with its full
original hypotheses. If \(C,Z\ge0\), \(\mu=\E Z\ge0\), and
\fitdisplay{
 C\le p+(6/5)\sqrt d-\kappa d,\qquad
 Z+\mu\le q+1+\sqrt d-(13/8+3q/8)d,\qquad
 \mu\le1-11d/8,
 \label{eq:q117-unified:strong-nine-data}
}{Q117.6.14}
then
\fitdisplay{
 C[(Z+\mu)^2+2\mu^2]\le9.
 \label{eq:q117-unified:strong-nine}
}{Q117.6.15}
This is Lemma~\ref{lem:reference-coefficient-nine}, case \(s\).
The actual S4 and S4F messages satisfy these hypotheses:
\(Z\le(1-\gamma)q+y-a_1\),
\(\gamma\ge(1-3v/2)d\), \(a_1\ge\kappa vd\), and
\(1-\mu\ge(2-3v/2)d\).
At \(v\le\theta\), the resulting coefficient of \(q\) is at
least \(25/64>3/8\), the other loss is at least
\(1671/1024>13/8\), and the mean loss is at least
\(89/64>11/8\).

\paragraph{The remaining point prerequisites.}
For S5, the exact data of the unit connector give
\(\gamma\ge d-3\delta_2/2>3d/8\) and
\(k|KR_1|\le(6/5)\sqrt{\theta d}<4y/5\).
Its mean is at most \(1-5d/8\), so
\(Z+\mu\le q+1+4y/5-d\).
For S5F, \(\gamma>7d/10\) and
\(k|KR_1|\le11y/12\); the latter follows from
\((36/25)\theta(1+\theta)<121/144\).
Together with \(\mu\le1-d\), this gives its point upper below.
For D2F and its two actual \(a=3\) prefixes,
\(\delta/(k\delta_2)\le39/160\) and
\(\|1-kB^3\|^2\le1+(36/25)\theta^3<(21/20)^2\).
The actual row formula therefore gives
\(Z\le q+21y/20\), with \(\mu\le1-d\).

For the D4 rows, use
\eqref{eq:q117-unified:long-square-moments}.
At \(b=3\), \(h\ge19/48\), \(\nu_0,N\le1\), whence
\[
 Z+\mu\le q+1+y-(19/48)(q+4)d\le q+22/19.
\]
The last inequality completes the square
\(y-19y^2/12\le3/19\).
At \(b=4\), \(\nu_0\le1+d/6\), and the coefficient of
\(q\) is at most one. The remaining upper is
\(1+y-11y^2/6+y^3+24y^4/5\), increasing on \([0,1/2]\)
and equal to \(22/15<3/2\) at \(1/2\).
For its derivative, use \(1-11y/3>0\) on \([0,1/4]\);
on \([1/4,1/2]\) the lower quadratic
\(1-11y/3+39y^2/5\) is positive.
For \(b\ge5\), the \(q\) coefficient is at most
\(1+d/10+6d^2/5\le267/250<107/100\).
The remaining \(Z\) term is at most
\((1+7d/5)\sqrt d-d/2+6d^2/5\), increasing and below
\(131/250<21/40\) at \(d=1/5\), using \(\sqrt{1/5}<9/20\).
In all three D4 rows \(\mu\le1\).

For low two-step \(b\ge5\), the actual finite row gives
\[
 Z+\mu\le q+1+10y/7+e(q)d,\qquad\mu\le1,
 \quad e(q)=\begin{cases}-1/10,&q\le1,\\6q/5-13/10,&q\ge1.
                         \end{cases}
\]
It uses \(\|\mathsf U_1\|\le10/7\),
\(\varepsilon\le3d/5\), and \(\delta\le6d/5\).
The \(b=4\) row has the stronger upper
\(q+1+27y/20-5d/6+\delta(q-1)\), with
\(\delta\le25d/24\). For \(q\le1\), the difference from
the preceding common upper is \(11y/140+11d/15\ge0\).
For \(q\ge1\), it is
\(11y/140+(19q/120+23/40)d\ge0\).
On the middle \(b\ge5\) domain, use
\(\varepsilon\le2d/3\), \(\delta\le4d/3\),
\(\|\mathsf U_1\|\le10/7\), and
\(\mu\le1-2d/3\le13/15\).
The two point uppers are \(q+12/7\) for \(q\le1\), and
\(4q/3+29/21\) for \(q\ge1\).
At \(b=4\), the middle mean is at most \(4/5\).
The stronger row just displayed is at most \(q+22/15\) when
\(q\le1\), since \(27y/20-5y^2/6\le7/15\) for
\(y\le1/2\). For \(q\ge1\), it is at most
\(19q/15+5/4\): the difference has nonnegative \(q\) coefficient
\(4/15-25y^2/24\ge1/160\), and at \(q=1\) is at least
\(31/60-27y/20+5y^2/6\ge1/20\).

For the three-step low \(b\ge5\) packet,
\(\ell\ge d\), \(\|\mathsf Q_2\|\le7/5\),
\(\varepsilon\le3d/5\), and \(\mu\le1-5d/4\) give
\[
 Z+\mu\le q+1+7y/5-21d/20\le q+22/15.
\]
On its middle domain use \(\|\mathsf Q_2\|\le3/2\),
\(\varepsilon\le2d/3\), and \(\mu\le3/4\); then
\(Z+\mu\le q+1+3y/2-11d/12\le q+31/20\).
At \(b=4\), the inputs \(\|\mathsf Q_2\|<10/7\),
\(\varepsilon\le d/6\), \(\ell\ge d\), and
\(\mu\le1-7d/4\) give
\(Z+\mu\le q+1+10y/7-29d/12\le q+1+3/14<q+5/4\).
Here the maximum of the last quadratic increment is
\(300/1421<3/14\).
At \(b=3\), \(\delta\le13/240\) gives
\(\|\mathsf Q_2\|^2\le87409/57600<25/16\).
The actual row gives \(Z\le(1-\eta)q-d+5y/4\);
using \(\mu\le1-d\),
\(Z+\mu\le q+1+5y/4-2d\le q+153/128<q+5/4\).

\paragraph{One rational certificate table.}
In Table~\ref{tab:q117-unified:curvature-data}, let
\(\mathcal P=M-\mathcal C(Y^2+2m_0^2)\).
The entries \(\mathcal C,Y,m_0\) are actual point and mean
uppers proved above. The last row is also used for the high-domain
bare path. Each row has \(\mathcal P\ge0\) on its full rectangle.
Here is an exact coefficient certificate: subdivide the indicated
\(q\) and \(y\) intervals into the stated equal grid, map each
cell affinely to \([0,1]^2\), and apply
\eqref{eq:bernstein-coefficients} with the printed bidegree.
Every Bernstein coefficient is nonnegative. By
Lemma~\ref{lem:bernstein-certificate}, this proves
\begin{equation}\label{eq:q117-unified:coefficient-bound}
 C_{a,t}[(Z+\mu)^2+2\mu^2]\le M
\end{equation}
for each actual row, including every boundary.
For the first three rows, \(0\le y\le26/45\); for the two
low two-step rows, \(0\le y\le9/20\).
The rows without \(y\) have degree zero in that variable.
Their grids and bidegrees determine every coefficient by a finite rational formula.

For the two low two-step entries abbreviate the already proved upper
by \(Y_\ell(q,y)=q+1+10y/7+e(q)y^2\).
The function \(K_*\) is defined with the row data above.
\begingroup\small\setlength{\tabcolsep}{3pt}
\begin{longtable}{@{}p{0.13\linewidth}p{0.10\linewidth}p{0.25\linewidth}p{0.10\linewidth}p{0.11\linewidth}p{0.10\linewidth}p{0.10\linewidth}@{}}
\caption{Fourth-coefficient certificates. A full \(q\) interval is
\([0,2]\); the two symbols \(-,+\) denote \([0,1],[1,2]\).
The grid lists the number of equal subdivisions in \((q,y)\).}
\label{tab:q117-unified:curvature-data}\\
\toprule
Rows & \(\mathcal C\)& \(Y\)&\(m_0\)&\(M\)&Degree&Grid\\
\midrule\endfirsthead
\toprule
Rows & \(\mathcal C\)& \(Y\)&\(m_0\)&\(M\)&Degree&Grid\\
\midrule\endhead
S5&\(\mathfrak C\)&\(q+1+4y/5-y^2\)&\(1-5y^2/8\)&\(37/4\)&\((3,6)\)&\((2,2)\)\\
S5F&\(\mathfrak C\)&\(q+1+11y/12-17y^2/10\)&\(1-y^2\)&10&\((3,6)\)&\((2,2)\)\\
D2F, D2F3L/H&\(\mathfrak C\)&\(q+1+21y/20-y^2\)&\(1-y^2\)&10&\((3,6)\)&\((2,2)\)\\
D4F&\(\mathfrak C_\infty\)&\(q+22/19\)&1&10&\((3,0)\)&\((8,1)\)\\
D4B&\(\mathfrak C_\infty\)&\(q+3/2\)&1&13&\((3,0)\)&\((2,1)\)\\
D4T&\(\mathfrak C_\ell\)&\(107q/100+61/40\)&1&13&\((3,0)\)&\((2,1)\)\\
D2BL/TL \(-\)&\(\mathfrak C\)&\(Y_\ell(q,y)\)&1&\(K_*(y^2)\)&\((3,6)\)&\((1,1)\)\\
D2BL/TL \(+\)&\(\mathfrak C\)&\(Y_\ell(q,y)\)&1&\(K_*(y^2)\)&\((3,6)\)&\((1,1)\)\\
D2TM \(-\)&\(\mathfrak C_\infty\)&\(q+12/7\)&\(13/15\)&14&\((3,0)\)&\((1,1)\)\\
D2TM \(+\)&\(\mathfrak C_\infty\)&\(4q/3+29/21\)&\(13/15\)&14&\((3,0)\)&\((1,1)\)\\
D2BM \(-\)&\(\mathfrak C_\infty\)&\(q+22/15\)&\(4/5\)&13&\((3,0)\)&\((1,1)\)\\
D2BM \(+\)&\(\mathfrak C_\infty\)&\(19q/15+5/4\)&\(4/5\)&13&\((3,0)\)&\((1,1)\)\\
D3TL&\(\mathfrak C_\ell\)&\(q+22/15\)&1&12&\((3,0)\)&\((4,1)\)\\
D3TM&\(\mathfrak C_\infty\)&\(q+31/20\)&\(3/4\)&12&\((3,0)\)&\((2,1)\)\\
D3F, D3B&\(\mathfrak C_\infty\)&\(q+5/4\)&1&11&\((3,0)\)&\((2,1)\)\\
Bare \(\bar f_2\)&\(\mathfrak C_\infty\)&\(q+5/4\)&\(3/4\)&\(75/8\)&\((3,0)\)&\((4,1)\)\\
\bottomrule
\end{longtable}
\endgroup

\begin{longtable}{@{}lp{0.36\linewidth}ccc@{}}
\caption{Integer scopes, continuous domains, and inverse payments.
Every row has \(a\ge3\) unless a smaller receiving range is printed.
S rows use \(qf_b\); D rows use \(h_2f_b\).}
\label{tab:q117-unified:scopes}\\
\toprule
Row&Integer range&\(d\)-interval&\(\sigma_r\)&\(\Delta_r\)\\
\midrule\endfirsthead
\toprule
Row&Integer range&\(d\)-interval&\(\sigma_r\)&\(\Delta_r\)\\
\midrule\endhead
S1&\(b\ge3,\ L\ge5\)&\([0,1/3]\)&\(\kappa\)&\(\frac{1}{600}\)\\
S2&\(b\ge3,\ L=4\)&\([0,1/3]\)&\(\kappa\)&\(\frac{1}{600}\)\\
S3L&\(b\ge3,\ L=3\)&\([0,1/4]\)&\(\kappa\)&\(\frac{1}{600}\)\\
S3H&\(b\ge3,\ L=3\)&\([1/4,1/3]\)&\(\kappa\)&\(\frac{1}{600}\)\\
S4&\(b\ge5,\ L=2\)&\([0,1/3]\)&\(1-v\)&\(\frac{3}{80}\)\\
S4F&\(b=3,4,\ L=2\)&\([0,1/3]\)&\(1\)&\(\frac{3}{80}\)\\
S5&\(b\ge4,\ L=1\)&\([0,1/3]\)&\(1-v\)&\(\frac{1}{600}\)\\
S5F&\(b=3,\ L=1\)&\([0,1/3]\)&\(1\)&\(\frac{1}{600}\)\\
S6&\(b=2,\ L\ge6\)&\([0,5/24]\)&\(0\)&\(\frac{3}{800}\)\\
D4F&\(b=3,\ L\ge4\)&\([0,1/3]\)&\(1\)&\(\frac{1}{80}\)\\
D4B&\(b=4,\ L\ge4\)&\([0,1/4]\)&\(2\)&\(\frac{1}{40}\)\\
D4T&\(b\ge5,\ L\ge4\)&\([0,1/5]\)&\(\frac{9}{4}\)&\(\frac{1}{50}\)\\
D2F&\(b=3,\ L=2,\ a\ge4\)&\([0,1/3]\)&\(1\)&\(\frac{1}{80}\)\\
D2F3L&\(b=3,\ L=2,\ a=3\)&\([0,1/6]\)&\(1\)&\(\frac{1}{80}\)\\
D2F3H&\(b=3,\ L=2,\ a=3\)&\([1/6,1/3]\)&\(1\)&\(\frac{1}{80}\)\\
D2BL&\(b=4,\ L=2\)&\([0,1/5]\)&\(2\)&\(\frac{1}{40}\)\\
D2BM&\(b=4,\ L=2\)&\([1/5,1/4]\)&\(2\)&\(\frac{1}{40}\)\\
D2TL&\(b\ge5,\ L=2\)&\([0,1/5]\)&\(\frac{9}{4}\)&\(\frac{1}{50}\)\\
D2TM&\(b\ge5,\ L=2\)&\([1/5,1/4]\)&\(2\)&\(\frac{1}{40}\)\\
D3F&\(b=3,\ L=3\)&\([0,1/3]\)&\(1\)&\(\frac{1}{80}\)\\
D3B&\(b=4,\ L=3\)&\([0,1/4]\)&\(2\)&\(\frac{1}{40}\)\\
D3TL&\(b\ge5,\ L=3\)&\([0,1/5]\)&\(\frac{9}{4}\)&\(\frac{1}{50}\)\\
D3TM&\(b\ge5,\ L=3\)&\([1/5,1/4]\)&\(2\)&\(\frac{1}{40}\)\\
\bottomrule
\end{longtable}
\begin{longtable}{@{}lcccc@{}}
\caption{Certificates for the two whole payment polynomials.
The bidegrees are those of \(P_r^+,P_r^0\) after the substitution
\(v=d+(\theta-d)t\); a degree-zero second variable is omitted.
Every coefficient in every listed cell is at least \(1/100\).}
\label{tab:q117-unified:scalar-certificates}\\
\toprule
Row&Degree \((+,0)\)&Cells \((c_+,c_0)\)&Coefficients&\(d\)-cut for \(c_+=2\)\\
\midrule\endfirsthead
\toprule
Row&Degree \((+,0)\)&Cells \((c_+,c_0)\)&Coefficients&\(d\)-cut for \(c_+=2\)\\
\midrule\endhead
S1&\(4,5\)&\(2,1\)&16&\(1/6\)\\
S2&\(4,5\)&\(1,1\)&11&\(-\)\\
S3L&\(5,5\)&\(1,1\)&12&\(-\)\\
S3H&\(5,5\)&\(1,1\)&12&\(-\)\\
S4&\((14,10),(10,5)\)&\(2,1\)&396&\(1/6\)\\
S5&\((9,5),(10,5)\)&\(1,1\)&126&\(-\)\\
S5F&\((6,2),(5,2)\)&\(1,1\)&39&\(-\)\\
S6&\((13,5),(5,3)\)&\(1,1\)&108&\(-\)\\
D4F&\(8,5\)&\(1,1\)&15&\(-\)\\
D4B&\(8,5\)&\(1,1\)&15&\(-\)\\
D4T&\(8,5\)&\(1,1\)&15&\(-\)\\
D2F&\((10,6),(5,2)\)&\(1,1\)&95&\(-\)\\
D2F3L&\((9,6),(5,2)\)&\(1,1\)&88&\(-\)\\
D2F3H&\(4,5\)&\(1,1\)&11&\(-\)\\
D2BL&\(4,5\)&\(1,1\)&11&\(-\)\\
D2BM&\(4,5\)&\(1,1\)&11&\(-\)\\
D2TL&\(4,5\)&\(2,1\)&16&\(1/10\)\\
D2TM&\(4,5\)&\(1,1\)&11&\(-\)\\
D3F&\((9,6),(5,1)\)&\(1,1\)&82&\(-\)\\
D3B&\(4,5\)&\(1,1\)&11&\(-\)\\
D3TL&\(4,5\)&\(1,1\)&11&\(-\)\\
D3TM&\(4,5\)&\(1,1\)&11&\(-\)\\
\bottomrule
\end{longtable}
S4F is paid by the S4 row and the exact stronger-prefix scale
comparison \eqref{eq:q117-unified:shortone-prefix-scale}.

\subsubsection{The common continuous payment}
\label{sec:q117-unified:payments}

Table~\ref{tab:q117-unified:scopes} specifies the complete integer
and continuous ranges of the parameter rows. A closed interval
starting at zero includes the polynomial extension at \(d=0\);
the actual zero-variance model is treated below. In each row the
finite data and point bounds just proved give
\begin{equation}\label{eq:q117-unified:two-actual-endpoints}
 \E[f_a\{T^LF\}^4]
 \le m^8\max\{M_rV_r,W_r\}.
\end{equation}
For the rows with a receiving-endpoint comparison, this means
the two actual endpoints followed by the proved comparison,
rather than a replacement of the covariance at an interior mean.

The outside scale in every row has
\(m\le1-\sigma_rd/2\), with \(\sigma_r\) in the scope table.
Define the finite inverse polynomial and its two exact payment
polynomials by
\begin{equation}\label{eq:q117-unified:payment-polynomials}
\begin{aligned}
 S_{5,r}&=\sum_{j=0}^{5}\binom{7+j}{7}(\sigma_rd/2)^j,\\
 P_r^+&=\frac{S_{5,r}-W_r}{d}-\Delta_r,\\
 P_r^0&=S_{5,r}-M_rV_r-\Delta_rd.
\end{aligned}
\end{equation}
The quotient is a polynomial, since both numerator terms have
constant term one in \(d\). The positive inverse-binomial series
gives \(m^{-8}\ge(1-\sigma_rd/2)^{-8}\ge S_{5,r}\).
All its denominators are positive on the displayed domains.

For each row depending on \(v\), set
\(v=d+(\theta-d)t\), \(0\le t\le1\).
Since \(\theta-d\ge7/96>0\), this parametrizes exactly
\(d\le v\le\theta\). For the other rows, the polynomial is
constant in \(t\). Table~\ref{tab:q117-unified:scalar-certificates}
gives a rational certificate for both entire polynomials in
\eqref{eq:q117-unified:payment-polynomials}.
The two bidegrees are in \((d,t)\); \(c_+,c_0\) give the number
of closed \(d\)-cells. Subdivide only at the indicated midpoint
when \(c_+=2\); in every cell the \(t\)-interval is \([0,1]\).
After the affine map, compute the Bernstein coefficients by
\eqref{eq:bernstein-coefficients}. Every coefficient is at least
\(1/100\). The rational coefficient tables include
their complete values and affine power coefficients.
Lemma~\ref{lem:bernstein-certificate} therefore proves both
payment polynomials nonnegative throughout their domains.
Consequently
\begin{equation}\label{eq:q117-unified:paid-endpoints}
 \max\{M_rV_r,W_r\}\le S_{5,r}-\Delta_rd
       \le m^{-8}-\Delta_rd.
\end{equation}
The finite-prefix comparison
\eqref{eq:q117-unified:shortone-prefix-scale} extends the S4
certificate to S4F, paying its entire additional nonzero endpoint.

For S1--S5F, \(m\ge5/6\) and \((5/6)^8>1/5\).
At S6, \(m=1\). For the short-two packets the respective
lower bounds on \(m^8\) are \(1/5\) at \(b=3\),
\(1/10\) at \(b=4\) or on the middle \(b\ge5\) domain,
and \(1/8\) on the low \(b\ge5\) domain.
Multiplying \eqref{eq:q117-unified:paid-endpoints} by the actual
\(m^8\) therefore gives the three required deficits on every
low and middle row. The S4 and S4F bound is stronger,
\begin{equation}\label{eq:q117-unified:shortone-two-step-strong}
 \E[f_a\{T^2(qf_b)\}^4]\le1-3d/400,
 \qquad a,b\ge3.
\end{equation}

\subsubsection{The high domains}
\label{sec:q117-unified:high-domains}

These remaining ranges use the same original terminal and
orientation. The needed path estimates are
\eqref{eq:compact-strong-path-means},
\eqref{eq:L56:claim}, and \eqref{eq:Q64:four}:
\[
 \E f_n=\E\bar f_n\le1-83d/32\quad(n\ge4),\qquad
 \E(\bar f_3^4/q)\le1-23d/36,\qquad
 \E\bar f_4^4\le1-d-4d^2.
\]
In particular \(\bar f_n\le2-5d\) for \(n\ge5\), and
\(\E f_a\le1-2d\) for \(a\ge3\).
From \(f_a\le5/2-q\),
\(f_aq\le25/16<8/5\).

For \(F=qf_2\), \(\E F=1-2d\). Positivity gives
\(T^LF\le2(1-2d)\bar f_{L-1}\).
At \(L\ge6,d\ge5/24\), its fourth moment is at most
\[
 (1-2d)[2(1-2d)(2-5d)]^4
 \le45927/49152<15/16<1-3d/800.
\]
The factors are positive and decreasing, so it suffices to use
their left endpoint; the displayed rational bound also uses
\(161/144<9/8\).

For \(F=h_2f_b,b\ge5,d\ge1/5\), the true terminal mean
from Section~\ref{proof:Q5-C073-connectors} is at most \(91/200\).
Thus \(T^LF\le(91/100)\bar f_{L-1}\).
At \(L=4\), use the inverse \(q\)-weighted \(\bar f_3\)
estimate and \(f_aq<8/5\); the resulting budget is below
\(35/36\). At \(L=5\), use \(f_a\le2\) and the
\(\bar f_4\) estimate, giving a budget below \(8/9\).
At every \(L\ge6\), \(\bar f_{L-1}\le1\), and the budget
is at most \((91/100)^4\). Each is below \(1-d/400\).
For \(b=4,d\ge1/4\), the same terminal calculation gives
\(TF\le4/5\); the identical three path ranges have common
budget \(512/625<1-d/400\).

For the two short high-domain connectors use the actual finite
terminal estimate \eqref{eq:L77:means}:
\(T(h_2f_b)\le13/18\) for every \(b\ge4,d\ge1/4\).
At \(L=2\), positivity gives \(T^2F\le(13/18)q\).
Since \(\E q^3=1+3d-\E g^3\le1+4d\le7/3\),
\begin{equation}\label{eq:q117-unified:high-two-step}
 \E[f_a(T^2F)^4]
 \le(13/18)^4(25/16)(7/3)
 =4998175/5038848<1-d/400.
\end{equation}

For \(L=3\), the actual bare \(\bar f_2=q+Kg\) has mean
\(1-d\), variance at most \(45d/32-d^2\), and covariance
\[
 \operatorname{Cov}(C_{a,0},\bar f_2)
 =-\int\operatorname{Re}\{G_a(1+z)\}\,d\nu-\beta_{a+1}d<0,
\]
since \(\operatorname{Re}\{G_a(1+z)\}\ge1-2x>0\).
Its point upper satisfies
\(\bar f_2+(1-d)\le q+1+\sqrt d-d\le q+5/4\),
and \(1-d\le3/4\). The last row of
Table~\ref{tab:q117-unified:curvature-data} gives coefficient
\(75/8\). The exact fourth identity with mass at most one and
nonpositive covariance yields
\begin{equation}\label{eq:q117-unified:bare-high-fourth}
 \E[f_a\bar f_2^4]
 \le(3/4)^4+(75/8)(103/288)=1409/384.
\end{equation}
Here \(45d/32-d^2\) increases on \([1/4,1/3]\) and is
at most \(103/288\). Therefore
\begin{equation}\label{eq:q117-unified:high-three-step}
 \E[f_a(T^3F)^4]
 \le(13/18)^4(1409/384)
 =40242449/40310784<1-d/400.
\end{equation}
The exact gaps in the last two comparisons are
\(40673/5038848\) and \(68335/40310784\), respectively;
both exceed \(1/1200\).

\paragraph{The degenerate model, decorations, and reversal.}
When \(d=0\), \(g=0\) on every positive-mass atom.
Successive leaf integration makes every actual arm and terminal
message identically one, even if \(K\ne0\). Thus each reference
moment is one, with no moment ratio required.
For \(d>0\), the low, middle, and high intervals above cover
their shared boundaries. Actual equal-pair products lie in
\([0,1]\); positivity of \(T\) restores \(E_t,E_o\) in
the three estimates. This proves
Proposition~\ref{prop:q117-unified:families}.

Whole reversal gives the same three bounds with receiver
\(\bar f_a\), propagation \(A^L\), and terminals
\(p\bar f_b\), \(p\bar f_2\), and \((Aq)\bar f_b\),
respectively. It reverses every actual edge at once.
With short terminal-arm length \(s=1\) or \(2\), an original
decorated component has \(a+b+L+s+2S_t+2S_o\) edges and
one more vertex. The fourth verification tree has
\(a+4(b+L+s)+8S_t+2S_o\) edges and one more vertex;
the receiving arm occurs once.

